\begin{abstract}
We develop four exact continuations of the ProveIt program on
polylogarithms, harmonic numbers, Gamma functions, and zeta derivatives.
First, for an arbitrary polynomial numerator in generalized harmonic
numbers, we classify the cancellation of every nonpositive even
spectral principal part at the centered shift $1/2$. The criterion is
invariance under the simultaneous reflection of all even harmonic
tails; its necessity follows from a theorem on formal difference
equations. We then evaluate every centered even moment of the square
of trigamma as a rational number plus a rational multiple of $\zeta(3)$,
and sum these values in a convergent polylogarithmic generator.

Second, a cotangent resolvent gives a completed formula for all
generalized Stieltjes indices and all argument derivatives. We
identify the entire local correction polynomial before coefficient
extraction. Partial fractions yield every rational cotangent kernel
bounded at infinity with no real pole. The same transform handles
continuous polygamma orders and all balanced Gamma primitives, with
an explicit conversion between analytic continuation in the order
and pointwise endpoint finite parts.

Third, fixed nonpositive indices can be eliminated at arbitrary
positions and depths in the entire relative harmonic interpolant.
We prove a canonical finite polynomial normal form and a genuinely
convergent Lerch generator for arbitrarily many zero indices surrounding
one free order. Classical polynomial Stieltjes primitives are recorded
with arbitrary endpoints and fully explicit coefficients.

Finally, every admissible quartic Tornheim ray is reduced to three
specified coordinates, and cyclic quartic combinations retain only
the diagonal coordinate. At all degrees we determine the exact
homogeneous freedom under the stated symmetry, axis, and Euler-plane
relations. An explicit deformation proves why diagonal data cease
to determine the cyclic fifth layer. The article includes complete
proofs, a source audit and integration notes, reproducible exact and numerical
checks, and a concrete further-research agenda.
\end{abstract}
